\theoremstyle{atheorem}
\newtheorem{claim}[theorem]{Claim}

\chapter{Pen-and-Paper State-Separating Proof}
\label{chp:penpaperssp}
This chapter introduces the notion of a State Separating Proof (SSP), and how it differs from a traditional code-based game-playing proof. Then, we define our pen-and-paper \acrlong{ssp} for the Free XOR \acrlong{gs} here.
\section{State Separating Proofs}
\label{sec:ssp}
\acrlong{ssp}s are a style of proofs built upon code-based game-playing proofs discussed by \citeauthor{EC:BelRog06} \cite{EC:BelRog06}. They define code-based game-playing as proofs that encompass an adversary's interaction with its environment using code (could be pseudocode or any other formalized programming language), and these interactions are referred to as a game. The game contains subroutines that initialize the state of the game, then run the adversary program. The adversary can call oracles, which are functions exposed by the game, and the game concludes by running a finalize subroutine which outputs a bit from the adversary. The adversary has a similar interaction between two different versions of the game, one in a real world and one in an ideal world, whose behaviors are parameterized by a bit $b=0$ and $b=1$, respectively. 
The success of the adversary is characterized by its `advantage', defined as the difference in the probability of the adversary returning the same bit in each game. In cryptographic notation, given two games $G_\mathit{real}$ and $G_\mathit{ideal}$, 
\[ \advantage{}{\adv, G^0_\mathit{real}, G^1_\mathit{ideal}} = \lvert \prob{A\to G_\mathit{real}=1} - \prob{A\to G_\mathit{ideal}=1} \rvert\]

Games and Oracles are constructs of a game-playing proof. 
\begin{itemize}
    \item \textbf{Oracles} are defined as a piece of code that provides some functionality. They usually either return a value, or modify the state of the Game.
    \item \textbf{Games} are a complete set of oracles that do not call other oracles, an initialize procedure and a finalize procedure that returns a bit. An adversary interacts with a game.
\end{itemize}
However, \acrlong{ssp}s differ by introducing the notion of a package. A package has its own local state, and it functions as an intermediate layer between a game and a package. Games contain packages, and packages contain oracles. In a \acrlong{ssp} \cite{brzuska_state_2018}, they are defined as follows:
\begin{itemize}
    \item \textbf{Oracles} are a piece of executable code. They usually either return a value or modify the state of the package they are contained in. Oracles in a package cannot call each other to prevent recursion.
    \item \textbf{Packages} are defined as a collection of related oracles with a shared state.
    \item \textbf{Games} are defined as a package with no dependencies on oracles of other packages $([G\to] = \varnothing)$.
    \item An \textbf{adversary} can be defined as a special package that is not a dependency of another package. It always outputs a single bit on termination $([\to A] = \varnothing)$.
\end{itemize}

\acrfull{ssp}s allow us to restrict the access to the state of each package to oracles in that package, and only provide the oracles themselves as an interface to other packages. This allows them to be highly modular and composable, allowing us to reuse a package used in one reduction in another, all while allowing changes in state that do not make changes to the whole game. If these changes cascaded to the whole game, it would require the cryptographer to once again verify the correctness of the game in the new state. \cite{brzuska_state_2018} formally defines the constructs we use here, and the mathematical properties of packages such as associativity and compositions. We slightly diverge from their notation of sequential composition by using the $\to$ operator instead of $\circ$. 


\section{On the Privacy of the Garbling of Various Gates}
\label{sec:sspANDgate}
In our work, we are concerned with the garbling of NOT, XOR and AND gates, since every boolean circuit can be represented with that combination of gates (functional completeness of boolean gates). Despite the fact that NOT and AND gates by themselves are functionally complete, we are interested in XOR gates since they are the basis on which the Free XOR optimization improves the efficiency of evaluating a \acrlong{gc}. We discuss these gates in order:
\subsection{NOT Gate}
\label{subsec:notgate}
\begin{figure}[ht]
    \centering
    \begin{minipage}[c][2cm][b]{0.30\linewidth}
        % \centering
        \begin{tabular}{cc}
            \toprule
            Bit-value & Wire label\\
            \midrule
            $0$ & $W$\\
            $1$ & $W \xor \Delta$\\
            \bottomrule
        \end{tabular}
        % \vspace{0.4cm}
        % (a) Mapping of wire labels and bit-values before a NOT gate
    \end{minipage}
    \begin{minipage}[c][2cm][b]{0.33\linewidth}
        % \centering
        \ctikzset{logic ports=ieee}
        \begin{tikzpicture}[transform shape]
            % Paths, nodes and wires:
            \node[ieeestd not port](N1) at (8.331, 7.47){} node[anchor=center] at (N1.text){};
            \node[shape=rectangle, minimum width=0.715cm, minimum height=0.465cm](N2) at (6.875, 7.5){} node[anchor=center] at (N2.text){$A$};
            \node[shape=rectangle, minimum width=0.715cm, minimum height=0.465cm](N4) at (9.75, 7.5){} node[anchor=center] at (N4.text){$A \xor \Delta$};
        \end{tikzpicture}
        % \vspace{0.4cm}
        % (b) AND gate
    \end{minipage}
    \begin{minipage}[c][2cm][b]{0.30\linewidth}
        % \centering
        \begin{tabular}{cc}
            \toprule
            Bit-value & Wire label\\
            \midrule
            $0$ & $W \xor \Delta$\\
            $1$ & $W$\\
            \bottomrule
        \end{tabular}
        % \vspace{0.4cm}
        % (c) Mapping of wire labels and bit-values after a NOT gate
    \end{minipage}
    \vspace{0.4cm}

    \caption{Wire label mappings before (left) and after (right) a NOT gate.}
    \label{fig:garbling-not-gate}
\end{figure}
\citeauthor{ALP:KolSch08} note that NOT gates do not need to be garbled separately at all \cite[Section 3.1]{ALP:KolSch08}. A NOT gate simply inverts the mapping between the bit-value of a wire and its wire label, i.e., if the wire label $W_b$ corresponds to bit $b$ before encountering a NOT gate, the wire label $W_b$ will correspond to bit $b\xor 1$ after the NOT gate (Figure \ref{fig:garbling-not-gate} describes this inversion). 
Since NOT gates are not garbled, they are not part of the security proof of the privacy of the garbling of a circuit.

\subsection{XOR Gate}
\label{subsec:xorgate}
Similarly to NOT gates, XOR gates are garbled for `free' in the Free XOR scheme. What this means in practice is that the \cevaluator directly computes the output wire label of an XOR gate using the input wire labels (as discussed in Section \ref{subsec:freexor}). There is no garbled table that the \cevaluator receives, and no decryption that takes place, and the security of the XOR gate depends directly on the secrecy of the circuit offset $\Delta$. As long as the adversary does not know $\Delta$, the method of computing XOR gates in Free XOR remains secure.

\subsection{AND Gate}
\label{subsec:andgate}

Since the privacy of the gates discussed so far is trivial, the privacy of the garbling of AND gates in the Free XOR scheme becomes the main focus of this work. As mentioned in Section \ref{subsec:freexor}, we base our \acrlong{ssp} on \citeauthor{JC:Applebaum16}'s proof in \cite{JC:Applebaum16}.

\begin{figure}[ht]
\setlength{\abovecaptionskip}{0pt}
\centering
    \procedureblock[]{$G^0_\mathsf{Sim\text{-}FreeXOR}$}{
        \underline{\underline{\mathsf{GARBLE\mathit{(f, x)}}}} \\
        c \gets f(x) \\
        (a, b) \gets x \\
        \Delta \sample \bin^\secpar \\
        A_a, B_b, C_c \sample \bin^\secpar \\
        A_{a\xor 1} \gets A_a \xor \Delta; B_{b\xor 1} \gets B_b \xor \Delta; C_{c\xor 1} \gets C_c \xor \Delta \\
        R_{00}, R_{01}, R_{10}, R_{11} \sample \bin^\secpar \\
        F: \text{Table} \\
        F[a][b]             \sample (\enc(A_a, R_{00}),          \enc(B_b, R_{00}           \xor C_{a\land b}))                 \\
        F[a][b\xor 1]       \sample (\enc(A_a, R_{01}),          \enc(B_{b\xor 1}, R_{01}   \xor C_{a\land({b\xor 1})}))          \\
        F[a\xor 1][b]       \sample (\enc(A_{a\xor 1}, R_{10}),  \enc(B_b, R_{10}           \xor C_{({a\xor 1})\land b}))         \\
        F[a\xor 1][b\xor 1] \sample (\enc(A_{a\xor 1}, R_{11}),  \enc(B_{b\xor 1}, R_{11}   \xor C_{({a\xor 1})\land ({b\xor 1})})) \\
        X \gets (A_a, B_b) \\
        d \gets  \{C_c: c\} \\
        \pcreturn (F,X,d) \\
       % 
    }

\caption{Code for $G^0_\mathsf{Sim\text{-}FreeXOR}$}
\label{fig:simfreexor-real-game}
\end{figure}


Since we use the simulation based definition of privacy (discussed in Section \ref{sec:defininigprivacy}), we define the garbling of an AND gate in the Free XOR scheme as the security game $G^0_\mathsf{Sim\text{-}FreeXOR}$ (defined in Figure \ref{fig:simfreexor-real-game}). The game exposes a $\mathsf{GARBLE}$ oracle that provides the core functionality of garbling an AND gate. It first obtains the evaluation of the circuit $f$ with the input $x$ to determine the output it has to garble. The input $x$ is partitioned into the bit-values for wire $A$ and $B$. Then, it samples a random offset value, and wire labels for the two active input wires and active output wire. It also samples four random values, one for each row of the garbled table. The order in which these values are used is irrelevant since they are not used across multiple rows of the garbled table, therefore, they can be used independently of the index of the garbled table, as long as they are not reused.
The garbled table is then generated as described in Section \ref{subsec:freexor}. The garbled inputs are the wire labels sampled corresponding to the active wire bit-values. The decoding information $d$ is set as a map from the active wire label for the output wire to the bit-value of the output wire.

\begin{figure}[ht]
\setlength{\abovecaptionskip}{0pt}
\centering
    \procedureblock[]{$G^{b}_{\mathsf{LIN-RK-KDM}}$}{
        \underline{\text{State:}} \\
        \Delta :\text{global circuit offset}\\
        \underline{\underline{\mathsf{ENC\mathit{({key}, {msg} ,{z}\in\bin)}}}} \\
        \pcif \Delta = \bot : \pccomment{If $\Delta$ is unset} \\
        \t \Delta \sample \bin^{\secpar} \pccomment{ Sample a fresh $\Delta$}\\
        \pcif b = 0 : \\
        \t {ct} \sample \enc({key} \xor \Delta,\ {msg} \xor {z}\cdot \Delta) \\
        \pcif b = 1 : \\
        \t {ct} \sample \enc({key} \xor \Delta,\ 0^{\lvert{msg}\rvert}) \\
        \pcreturn {ct} \\
    }
\caption{Code for $G^{b}_{\mathsf{LIN-RK-KDM}}$}
\label{fig:linrkkdm-game}
\end{figure}


\citeauthor{JC:Applebaum16} defines the notion of \acrshort{linrkkdm} security of a symmetric encryption scheme $(\enc,\dec)$ as the computational indistinguishability of the two games $\mathsf{Real}_s$ and $\mathsf{Fake}_s$ \cite[Section 3.2]{JC:Applebaum16}, where 
\[\mathsf{Real}_s: (\Delta, M, b) \mapsto \enc_{S+\Delta}(M+bS) \text{, and,}\]
\[\mathsf{Fake}_s: (\Delta, M, b) \mapsto \enc_{S+\Delta}(0^{l\times N}) \text{.}\]

We redefine the same security notion to better fit our \acrshort{ssp} setting: An encryption scheme is said to be \acrshort{linrkkdm} secure if any adversary $\adv$ interacting with $G^{b}_{\mathsf{LIN-RK-KDM}}$ (Figure \ref{fig:linrkkdm-game}) in the real and ideal world cannot distinguish between them with more than negligible advantage, i.e.,
\[ \advantage{}{\adv, G^{0}_{\mathsf{LIN-RK-KDM}}, G^{1}_{\mathsf{LIN-RK-KDM}}} \leq \negl \text{.}\]

Our $G^{b}_{\mathsf{LIN-RK-KDM}}$ provides an interface $\mathsf{ENC}$ that takes the key, message and an extra parameter $z$ as input, and returns an encryption of the message in the real game, and an encryption of zeroes in the ideal game (Figure \ref{fig:linrkkdm-game}). It is important to note that the real \acrshort{linrkkdm} game does not simply return an encryption under the key used as input to the $\mathsf{ENC}$ oracle. It XORs this key by an offset of $\Delta$, and optionally XORs the same offset to the message, depending on the value of $z$.  

\newacronym{ppt}{PPT}{probabilistic, polynomial-time}
\begin{theorem}
    If $\enc$ is the \acrshort{linrkkdm} secure encryption scheme used in the Free XOR garbling scheme, then for f=\{AND\} and $\forall x\in \bin^2$, there exists a \acrfull{ppt} simulator $G^1_\mathsf{Sim\text{-}FreeXOR}$ and reduction $\rdv$ such that for all \acrshort{ppt} adversaries $\adv$,
    \[\mathsf{Adv}(\adv, G^0_\mathsf{Sim\text{-}FreeXOR}, G^1_\mathsf{Sim\text{-}FreeXOR}) \leq \mathsf{Adv}(\adv\to\rdv, G^{0}_{\mathsf{LIN-RK-KDM}}, G^{1}_{\mathsf{LIN-RK-KDM}}) \] 
\end{theorem}

\begin{figure}[ht]
\centering
    \procedureblock[]{$G^1_\mathsf{Sim\text{-}FreeXOR}$}{
        \underline{\underline{\mathsf{GARBLE\mathit{(f(x), \Phi(f))}}}} \\
        c \gets f(x) \\
        \Delta \sample \bin^\secpar //\text{global offset}\\
        A_0, B_0, C_0 \sample \bin^\secpar \\
        A_1 \gets A_0 \xor \Delta \\
        B_1 \gets B_0 \xor \Delta \\
        C_1 \gets C_0 \xor \Delta \\
        R_{00}, R_{01}, R_{10}, R_{11} \sample \bin^\secpar \\
        F: \text{Table} \\
        F[0][0] \sample (\enc(A_0, R_{0,0}),\enc(B_0, R_{0,0} \xor C_0)) \\
        F[0][1] \sample (\enc(A_0, R_{0,1}),\enc(B_1, R_{0,1} \xor C_1)) \\
        F[1][0] \sample (\enc(A_1, R_{1,0}),\enc(B_0, R_{1,0} \xor C_1)) \\
        F[1][1] \sample (\enc(A_1, R_{1,1}),\enc(B_1, R_{1,1} \xor C_1)) \\
        X \gets (A_0, B_0) \\
        d \gets  {C_0: 0} \\
        \pcreturn (F,X,d) \\
       % 
    }

\caption{Code for $G^1_\mathsf{Sim\text{-}FreeXOR}$}
\label{fig:simfreexor-ideal-game}
\end{figure}


We now describe the simulation strategy used by $G^1_\mathsf{Sim\text{-}FreeXOR}$ (Figure \ref{fig:simfreexor-ideal-game}). The simulator receives as input the output of the circuit evaluation, $f(x)$, which corresponds to the bit-value of the output wire of the AND gate. It then samples random bitstrings for the wire labels corresponding to the zero bit-values on the input and output wires. It also samples the global circuit offset $\Delta$, and random values for each row of the garbled table. 
The correctness property of garbling requires that when the \cevaluator evaluates a garbled table, they are able to decrypt exactly one row of the table with the input wire labels they have. It also requires that the final decrypted output wire label corresponds to the correct function evaluation $f(x)$.

Since the simulator knows $f(x)$, the simulation strategy involves  mapping one of the output wire labels (our simulator in Figure \ref{fig:simfreexor-ideal-game} chooses $C_0$) to the bit-value $f(x)$, and making sure that row that the adversary can decrypt actually encrypts the chosen output wire label. We also need to ensure that if the adversary receives wire labels $A_0$ and $B_0$ as the garbled input $X$, then every encryption on the table that uses these as the key is encrypted with real messages, but for the other encryptions, we can forge them by encrypting zeroes. This concludes the generation of the garbled table $F$. The garbled input $X$ is a tuple of the active wire labels of the first and second wire, which we assumed to be $A_0$ and $B_0$ in our simulator. The choice only matters to select which row of the garbled table is available to the adversary for decryption. Finally, the decoding information $d$ is a map from our chosen active output wire label that was encrypted with the active input wire labels, to the real value of $f(x)$ that is available to the adversary. Finally, the adversary returns this $(F,X,d)$ tuple.

\begin{figure}[ht]
\centering
\begin{minipage}[t]{0.45\linewidth}
    \procedureblock[]{$G^0_{\mathsf{Wrapper},{Gbl}}$}{
        \underline{\underline{\mathsf{GBL}(f,x)}} \\
        (F, X, d) \sample \mathsf{GARBLE}(f, x) \\
        \pcreturn (F, X, d) \\
    }
\end{minipage}
\begin{minipage}[t]{0.45\linewidth}
    \procedureblock[]{$G^1_{\mathsf{Wrapper},{Sim}}$}{
        \underline{\underline{\mathsf{GBL}(f,x)}} \\
        (F, X, d) \sample \mathsf{GARBLE}(f(x), \Phi(f)) \\
        \pcreturn (F, X, d) \\
    }
\end{minipage}

\caption{Code of $G^{b}_{wrapper}$ }
\label{fig:wrapper-game}
\end{figure}


Another problem we quickly realise is that the two $\mathsf{GARBLE}$ oracles expect different inputs. $\mathsf{GARBLE}$ in the real game takes a circuit $f$, which in this case is a simple AND operation, and it also takes the inputs to the circuit, $x$. The $\mathsf{GARBLE}$ oracle in the ideal game however, takes $f(x)$ and $\Phi(f)$. To solve this incompatibility, we provide another game $G^{b}_{wrapper}$ that provides a shared interface `$GBL$' that can be composed with the real and ideal $G^b_\mathsf{Sim\text{-}FreeXOR}$ games. However, we do not include this in our proof in order to keep it simple to understand and translate into Domino code.

\begin{lemma}
    Let $\rdv$ be the reduction defined in Figure \ref{fig:reduction}. Then, for all \acrshort{ppt} adversaries $\adv$,
    \[\mathsf{Adv}(\adv, G^0_\mathsf{Sim\text{-}FreeXOR}, G^1_\mathsf{Sim\text{-}FreeXOR}) = \mathsf{Adv}(\adv \to \rdv, G^{0}_{\mathsf{LIN-RK-KDM}}, G^{1}_{\mathsf{LIN-RK-KDM}}) \]
\end{lemma}

\begin{figure}[ht]
\setlength{\abovecaptionskip}{0pt}
\centering
    \procedureblock[]{$\rdv$}{
        \underline{\underline{\mathsf{GARBLE\mathit{(f, x)}}}} \\
        c \gets f(x) \\
        (a, b) \gets x \\
        A_a, B_b, C_c \sample \bin^\secpar \\
        R_{00}, R_{01}, R_{10}, R_{11} \sample \bin^\secpar \\
        F: \text{Table} \\
        F[a][b]             \sample (\enc(A_a, R_{00}),          \enc(B_b, R_{00}           \xor C_{c}))                 \\
        F[a][b\xor 1]       \sample (\enc(A_a, R_{01}),          \mathsf{ENC}(B_{b}, R_{01}\xor C_{c}, {a\land (b\xor 1)}))          \\
        F[a\xor 1][b]       \sample (\mathsf{ENC}(A_{a}, R_{10} \xor C_{c}, {(a\xor 1)\land(b)}),  \enc(B_b, R_{10}))         \\
        F[a\xor 1][b\xor 1] \sample (\mathsf{ENC}(A_{a}, R_{11}, {(a\xor 1)\land(b \xor 1)}),  \mathsf{ENC}(B_{b}, R_{11}\xor C_{c}, {(a\xor 1)\land(b\xor 1)})) \\
        X \gets (A_a, B_b) \\
        d \gets  \{C_c: c\} \\
        \pcreturn (F,X,d) \\
       % 
    }

\caption{Code for Reduction $\rdv$}
\label{fig:reduction}
\end{figure}


We now describe the reduction $\rdv$. The interface to the reduction is a $\mathsf{GARBLE}$ oracle. It initially computes $f(x)$ and samples three bitstrings $A_a$, $B_b$ and $C_c$ of length $\secpar$ assuming $a$,$b$ and $c$ are the active wire labels. It also samples four random bitstrings of length $\secpar$, just like the real $G^0_\mathsf{Sim\text{-}FreeXOR}$. However, unlike that game, the reduction does not sample its own $\Delta$. This is because $\Delta$ is part of the state of the \acrshort{linrkkdm} game, and therefore, the reduction must rely on $\mathsf{ENC}$ oracle access from the \acrshort{linrkkdm} game in order to generate any of the ciphertexts that require inactive wire labels. The rest of the reduction follows the same steps as $G^0_\mathsf{Sim\text{-}FreeXOR}$, and inlining the $\mathsf{ENC}$ oracle gives us code that is equivalent to $G^0_\mathsf{Sim\text{-}FreeXOR}$, except for one important difference. For the calculation of $F[a\xor 1][b]$, $\rdv$ encrypts the output wire label in the first ciphertext, as opposed to the second one which is where $G^0_\mathsf{Sim\text{-}FreeXOR}$ encrypts it.

\textcolor{red}{wip - the argument to show that }
\[F[a\xor 1][b]\sample (\mathsf{ENC}(A_{a}, R_{10} \xor C_{c}, {(a\xor 1)\land(b)}),  \enc(B_b, R_{10}))\]
\[\implies F[a\xor 1][b]\sample (Enc(A_{a\xor 1}, R_{10} \xor C_{c} \xor S*(a\xor1)\land(b)),  \enc(B_b, R_{10}))\], 
\[\implies F[a\xor 1][b]\sample (Enc(A_{a\xor 1}, R_{10} \xor C_{c\xor((a\xor 1) \land b)}),  \enc(B_b, R_{10}))\], 

and 
\[        F[a\xor 1][b]       \sample (\enc(A_{a\xor 1}, R_{10}),  \enc(B_b, R_{10}           \xor C_{({a\xor 1})\land b}))         \]
are identically distributed.

TODO:
\begin{itemize}
    \item Write about the code equivalence of $R->G_LINRKKDM-1 = G-SimFreeXOR1$
    \item List all the game hops that were discussed above and the contradiction argument on the reduction
    \item Done for this section
\end{itemize}